% Pursuing Stacks (cotes 134-2 to 134-8): archive cote 134-5, site page 58 (PDF page 59), typescript page 245.
% typescript. Draft compared: ps4.tex:2928-2996.
% AI-transcribed by Claude Opus 5.5 (claude-opus-5-5, adaptive thinking, effort high), Message Batches API, batch ps2 (msgbatch_01NVwr1Yuo8qJf5rykeGBXPK), 2026-11-12; not yet checked by eye.
\scanpage{134-5/58}
\mspage{245}

weak test category, namely that the functor
\[ i^{*}_{A} = j_{A} : (\mathrm{Cat}) \longrightarrow A^{\wedge} \]
be model-preserving (which implies that $i_{A}$ is too).

Following def.6 (p.176), w'll say that $i$ is a \uline{test} \uline{functor} if $i$
\uline{and} the induced functors $i_{/a} : A_{/a} \longrightarrow A \longrightarrow M$ are weak test functors.
Using the definition of a test category (def.3, p.173), we see that this
just means that the following conditions hold:

1') $A$ is a test category.

2') $(M,M_{as})$ is modelizing (i.e. same as 2) above).

3') The functor $i$ and the induced functors $i_{/a} : A_{/a} \rightarrow M$ are \xout{xxw}
$M_{as}$ - aspheric.

This last condition merits a name by itself, independently of the
other assumptions:

\uline{Definition} 2. Let $(M,M_{as})$ be any asphericity structure, $A$ a small cate-%
gory, and $i: A \rightarrow M$ a functor. We'll say that $i$ is \uline{totally} $M_{as}$-\uline{aspheric}
(or simply \uline{totally} \uline{aspheric} , if no confusion is feared), if $i$ and \xout{th}
the induced functors $i_{a} : A_{/a} \longrightarrow M$ are $M_{as}$-aspheric (for any $a$ in $A$),
We'll say that $i$ is \uline{locally} $M_{as}$-\uline{aspheric} (or simply \uline{locally} \uline{aspheric})
if for any $a$ in $A$, the induced functor $i_{/a} : A_{/a} \longrightarrow M$ is $M_{as}$-aspheric.

Thus $i$ is totally aspheric iff it is aspheric and locally aspheric.
On the other hand, $i$ is a test functor iff $A$ is a test category, $(M,M_{as})$
is modelizing, and $i$ is totally aspheric.

In order for $i$ to be locally aspheric, it is ns that \inserted[typed]{$i$ factor through $M_{as}$ and} for any
$x$ in $M_{as}$ , $i^{*}(x)$ be \uline{aspheric over} $e_{A^{\wedge}}$ ; if $B$ is any \xout{generating} subca-%
tegory of $M$ containing $i(A)$ and generating the asphericity structure,
it is enough in this latter condition to take $x$ in $B$ .

\inserted{\uline{Remarks.}}
Thus we see that the three gradations for the test-functor notion,
as suggested by the definitions 5 and 6 of section 65 and now by the
present context of asphericity structures, just amount to gradations \xout{on}
for the test conditions on the category $A$ itself (\inserted{namely} to be \inserted{a} pseudo-test, a
weak test or just a plain test -category), and a two-step gradation on
the asphericity condition for $i$ (namely, that $i$ be either just aspheric
in the two first cases, or totally aspheric in the third), \inserted{while} these two
asphericity conditions on $i$ \deleted{being} \inserted{are} of significance, \xout{xxxx} independently of
any specific assumption which we may make on either $A$ or $(M,M_{as})$ . This

\margin{58}
